\section{System Architecture}
% **Content:**
%   - Provide a high-level overview of the entire system.
%   - Include a main architectural diagram showing the two primary components:
%     1. The containerized web service for text analysis (backend).
%     2. The ontology editor instance with the custom extension (frontend).
%   - Illustrate the interaction flow: user inputs text in the ontology editor -> request to extension's backend -> request to text analysis service -> service processes text -> structured output returned -> output parsed and integrated into the ontology.
%   - Discuss key design decisions for this architecture (e.g., microservice approach for the pipeline, choice of REST API for communication).

The system presented in this thesis is designed as a cohesive framework that bridges the gap between unstructured clinical narratives and formal ontological representations. It comprises two primary, yet interconnected, components: a backend clinical text analysis service named onc2ont, and a frontend integration within WebProtégé, realized as the \textit{WebProtégé CTakes Extension}. This architecture facilitates the automated conversion of clinical text into structured OWL2 entities and their seamless incorporation into active ontology projects.

\subsection{High-Level System Overview}
The overall architecture promotes a decoupled design, where the intensive text processing tasks are handled by a specialized, containerized web service, while the user interaction and ontology integration are managed within the familiar WebProtégé environment. The onc2ont service acts as a powerful backend engine, performing Natural Language Processing (NLP) and semantic conversion. The WebProtégé CTakes Extension provides the user interface and the necessary \textit{glue} logic to invoke the onc2ont service and integrate its results.

% Suggestion: Place the simplified complete system overview diagram here.
% This diagram should show the user interacting with WebProtégé, the extension making a call to the onc2ont service, and the data flow.
% Example elements: User -> WebProtégé UI (CTakes Extension Portlet) -> WebProtégé Server (Extension Backend) --HTTP POST Request--> onc2ont Service (Docker Container) --OWL/TTL Response--> WebProtégé Server (Extension Backend) -> Ontology Update.
\begin{figure}[h!]
    \centering
    \includegraphics[width=0.9\textwidth]{MemoriaFinal/images/overall_system_simplified_v1.1.png} % Replace with your actual diagram file
    % \fbox{\parbox[c][10cm][c]{0.9\textwidth}{\centering PLACEHOLDER FOR OVERALL SYSTEM DIAGRAM HERE \\ \textit{Illustrating interaction between WebProtégé with CTakes Extension and the onc2ont service.}}}
    \caption{High-Level System Architecture Overview.}
    \label{fig:overall_system_architecture}
\end{figure}



The typical interaction flow through the system is as follows:
\begin{enumerate}
    \item \textbf{User Input:} The user navigates to the \textit{Clinical Text Converter} perspective within WebProtégé and inputs clinical text into the \textit{onc2ont Portlet}.
    \item \textbf{Client-Side Request:} The \textit{Onc2OntPortletPresenter} (GWT client-side component) captures the text and initiates a request to the WebProtégé server-side backend.
    \item \textbf{Server-Side Handling (WebProtégé):} The \textit{ProcessClinicalNotesActionHandler} on the WebProtégé server receives this request. It is responsible for communicating with the external onc2ont service.
    \item \textbf{External Service Invocation:} The \textit{ProcessClinicalNotesActionHandler} constructs an HTTP POST request containing the clinical text and sends it to the \textit{/process} endpoint of the onc2ont service.
    \item \textbf{Text Analysis and Semantic Conversion (onc2ont):} The onc2ont service, running in a Docker container, processes the text using its cTAKES-based pipeline and \textit{MyOwlCasConsumer} to generate OWL/TTL formatted semantic data.
    \item \textbf{Response Reception:} The onc2ont service returns the generated OWL/TTL data in the HTTP response to the \textit{ProcessClinicalNotesActionHandler}.
    \item \textbf{Ontology Integration:} The \textit{ProcessClinicalNotesActionHandler} parses the received TTL data using WebProtégé's OWL API with Rio Turtle format and utilizes WebProtégé's change management system to apply the new axioms and annotations to the user's currently active ontology project.
    \item \textbf{User Feedback:} The WebProtégé interface updates to reflect the newly integrated entities and provides processing statistics to the user.
\end{enumerate}

\subsection{Component Architectures}

\subsubsection{onc2ont: Clinical Text Analysis Service}
The onc2ont service is a self-contained Docker-containerized application that transforms unstructured clinical text into structured OWL/TTL semantic data. The service integrates the Apache cTAKES NLP pipeline with custom components (\textit{MyExtendingFileTreeReader}, \textit{MyOwlCasConsumer}) and provides a RESTful API endpoint for web-based clinical text processing.

The service architecture consists of three main layers:

\textbf{API Layer:} A Flask-based HTTP server exposes the \texttt{/process} endpoint, handling clinical text input and returning TTL-formatted semantic data. File I/O management and Java subprocess orchestration are implemented here.

% Implementation details: Flask API in api.py with subprocess execution of Java pipeline

\textbf{Processing Pipeline:} The core NLP processing is handled by the \textit{PlaintextToTtl} Java class, which orchestrates a complete cTAKES analysis chain including sentence detection, POS tagging, UMLS dictionary lookup, assertion analysis, and relation extraction.

% Complete pipeline implementation in PlaintextToTtl.java with cTAKES AnalysisEngine components

\textbf{Semantic Conversion:} The \textit{MyOwlCasConsumer} component transforms linguistic annotations into formal OWL ontology structures using Apache Jena, creating individuals, properties, and relations that represent medical concepts and their relationships.

% OWL model generation code in MyOwlCasConsumer.java using Apache Jena framework

\begin{figure}[hh]
    \centering
    \includegraphics[width=0.7\textwidth]{MemoriaFinal/images/onc2ont_service_architecture_simplified.png} 
    \caption{onc2ont Service Architecture.}
    \label{fig:onc2ont_architecture}
\end{figure}

\subsubsection{WebProtégé CTakes Extension}
The WebProtégé CTakes Extension provides seamless integration of clinical text processing capabilities within the ontology development environment. The extension follows WebProtégé's established client-server architecture pattern:

\textbf{Client-Side Components (GWT):} 
\begin{itemize}
    \item UI presentation layer with \textit{Onc2OntView} and \textit{Onc2OntViewImpl} components
    \item Business logic handled by \textit{Onc2OntPortletPresenter} for user interaction management
    \item Integration within the \textit{Clinical Text Converter} perspective as an \textit{onc2ont Portlet}
\end{itemize}

% GWT client implementation follows WebProtégé's portlet architecture patterns

\textbf{Server-Side Components (Java/Spring):}
\begin{itemize}
    \item \textit{ProcessClinicalNotesActionHandler} manages HTTP communication with the external onc2ont service
    \item Integration with WebProtégé's change management system for ontology modification
    \item TTL data parsing and OWL axiom generation using WebProtégé's OWL API integration
\end{itemize}

% Server-side action handler implementation using WebProtégé's command pattern and Spring framework

This extension leverages WebProtégé's existing infrastructure for portlet creation, perspective management, and server-side action handling, ensuring consistent integration with the platform's architecture.

% Suggestion: Place the simplified WebProtégé CTakes Extension diagram here.
% This diagram should show the client-server interaction within WebProtégé and its call to the external onc2ont service.
% Example elements: User Interface (onc2ontPortlet GWT) -> Onc2OntPortletPresenter (GWT) --RPC--> ProcessClinicalNotesActionHandler (WebProtégé Server) --HTTP--> External onc2ont Service. Response handling and ontology update also shown on server side.
\begin{figure}[h!]
    \centering
    \includegraphics[width=1.0\textwidth]{MemoriaFinal/images/webprotege_ctakes_simplified_v1.1.png} % Replace with your actual diagram file
    %\fbox{\parbox[c][9cm][c]{0.8\textwidth}{\centering PLACEHOLDER FOR SIMPLIFIED WEBPROTÉGÉ CTAKES EXTENSION DIAGRAM HERE \\ \textit{Illustrating client-server components of the WebProtégé extension and its interaction with the external service.}}}
    \caption{WebProtégé CTakes Extension Architecture.}
    \label{fig:webprotege_extension_architecture}
\end{figure}

\subsection{Key Design Decisions}
Several key design decisions underpin this architecture:
\begin{itemize}
    \item \textbf{Microservice Approach for onc2ont:} The text analysis pipeline is developed as an independent service. This promotes modularity, allowing the onc2ont service to be developed, deployed, and scaled independently of WebProtégé. It also allows other applications to potentially consume this service.
    \item \textbf{REST API for Inter-Service Communication:} Communication between the WebProtégé extension's backend and the onc2ont service is achieved via a RESTful HTTP API. This is a standard, language-agnostic, and stateless communication protocol, simplifying integration and ensuring interoperability.
    \item \textbf{Containerization of onc2ont:} Using Docker for the onc2ont service simplifies deployment, manages dependencies, and ensures a consistent runtime environment across different setups.
    \item \textbf{Leveraging WebProtégé's Extensibility:} The frontend is built as an extension to WebProtégé, utilizing its established mechanisms for UI components (portlets, perspectives) and server-side logic (actions, change management). This reduces boilerplate development and ensures tight integration with the ontology editing environment.
    \item \textbf{Decoupled Processing and Integration:} The tasks of heavy NLP processing (onc2ont) and ontology integration (WebProtégé extension) are clearly separated. This separation of concerns improves maintainability and allows for specialized optimization of each component.
\end{itemize}
This architectural design aims to provide a robust, scalable, and user-friendly system for integrating advanced clinical text analysis capabilities into a collaborative ontology engineering platform.

